% ===================================================================
%  সারণি নং 5 — বাড়ি ও তার নির্মাণ।
%
%  Ceci n'est PAS une transcription : c'est une traduction, faite en
%  2026, en bengali d'aujourd'hui. Voir texto/bn/10-sarani-01.tex
%  pour la consigne, qui vaut pour les seize fichiers.
%
%  LES DEUX NOMS SUIVENT LE HINDI, pour que les deux colonnes
%  indiennes s'accordent : M. Leblanc est « শ্রী শ্বেত », M. Roux
%  garde son nom translittere, « শ্রী রু ».
% ===================================================================

%%K t05-apar-1 apar
\VUsaut{6.0mm}
\VUtitre{15.0pt}{60}{সারণি নং 5}
\VUsaut{1.4mm}
\VUtitre{11.4pt}{0}{বাড়ি ও তার নির্মাণ।}
\VUsaut{2.2mm}

%%K t05-tit-1 sub
\VUpk{10.2pt}{40}{ভালো সাক্ষাৎ।}
\VUsaut{3.0mm}

%%K t05-01-1 p
\textsc{জন}। --- কাকা, আমি খুব জানতে চাই \VUgras{বাড়ি} কীভাবে তৈরি হয়।

%%K t05-01-2 p
\textsc{কাকা} \textsuperscript{(80)}। --- এ তো সহজ, বাছা; আমার সঙ্গে এসো, আমি
তোমাকে সেই বাড়িটি দেখাব যা শ্রী বার্নার্ড \textsuperscript{(79)} বানাচ্ছেন
এবং যেখানে আমি থাকতে যাচ্ছি।

%%K t05-01-3 p
\textsc{জন}। --- আর এই বাড়ির \VUgras{নকশা} \textsuperscript{(76)} কে
এঁকেছেন?

%%K t05-01-4 p
\textsc{কাকা}। --- \VUgras{স্থপতি} \textsuperscript{(75)}; তিনি খুব স্পষ্ট
একটি ছবি এঁকেছিলেন, তা মঞ্জুর হল, আর \VUgras{ঠিকাদার} \textsuperscript{(77)}
কাজ শুরু করে দিলেন।

%%K t05-01-5 p
\textsc{জন}। --- আর ঠিকাদার কে?

%%K t05-01-6 p
\textsc{কাকা}। --- শ্রী শ্বেত, তোমার বন্ধু জেমসের বাবা। তিনি স্থপতি শ্রী রু-র
সঙ্গে মিলে মজুরদের নির্দেশ দেন।

%%K t05-01-7 p
এই দেখো! শ্রী শ্বেত ঠিক সময়েই নিজের \VUgras{মাপকাঠি} \textsuperscript{(78)}
নিয়ে আসছেন, আর শ্রী রু নিজের নকশা নিয়ে।

%%K t05-01-8 p
নমস্কার, মহাশয়েরা। আপনারা কেমন আছেন?

%%K t05-01-9 p
\textsc{শ্রী রু}। --- খুব ভালো, ধন্যবাদ। নমস্কার জন; ছেলেটি কত বড় হয়েছে!

%%K t05-01-10 p
\textsc{কাকা}। --- সত্যিই, ও পুরো একজন ছোট মানুষ। ও খুব গম্ভীর; যা দেখে তারই
ব্যাখ্যা করতে হয়। আজ ও জানতে চায় বাড়ি কীভাবে তৈরি হয়।

%%K t05-tit-2 sub
\VUpk{10.2pt}{40}{মজুরেরা।}
\VUsaut{3.0mm}

%%K t05-01-11 p
\textsc{শ্রী শ্বেত}। --- তবে আমার সঙ্গে চলো, আমার ছোট্ট বন্ধু, এখনই বুঝিয়ে
দিচ্ছি, কারণ শিখতে চাওয়া ছেলেমেয়েদের আমি খুব ভালোবাসি।

%%K t05-01-12 p
ওই মজুরটিকে দেখছ, যার হাতে \VUgras{বেলচা} \textsuperscript{(82)}: সে
\VUgras{খননকারী} \textsuperscript{(81)}। জলের নল বসানোর জন্য সে একটি
\VUgras{খাদ} \textsuperscript{(46)} খুঁড়ছে। সে-ই সেই মাটিও খুঁড়েছিল যেখানে
বাড়িটি দাঁড়িয়ে, যাতে রাজমিস্ত্রি চারটি \VUgras{দেয়াল} তুলতে পারে। সেই
দেয়ালের যে অংশ মাটির নিচে তাকে বলে \VUgras{ভিত} \textsuperscript{(2)}। ভিতে
ঘেরা জায়গাটি \VUgras{ভাঁড়ারঘর} \textsuperscript{(3)} গড়ে। এই ভাঁড়ারঘরে একটি
ছোট জানালা আছে যাকে বলে \VUgras{ঘুলঘুলি} \textsuperscript{(5)}, একটি
\VUgras{দরজা} \textsuperscript{(4)} ও একটি সিঁড়ি।

%%K t05-01-13 p
\VUgras{রাজমিস্ত্রি} \textsuperscript{(108)} দেয়াল তুলে যেতে লাগল,
\VUgras{দরজার} \textsuperscript{(8)} ও \VUgras{জানালার} \textsuperscript{(17)}
জন্য জায়গা ছেড়ে।

%%K t05-01-14 p
\textsc{জন}। --- কিন্তু সেই রাজমিস্ত্রিকে তো আমি দেখতেই পাচ্ছি না!

%%K t05-01-15 p
\textsc{শ্রী শ্বেত}। --- বাড়ির কাজ তার শেষ হয়ে গেছে। সে ওই
\VUgras{আস্তাবলের} \textsuperscript{(54)} পাশে, নিজের \VUgras{ভারায়}
\textsuperscript{(110)} উঠে, \VUgras{ধোপাঘরের} \textsuperscript{(56)} দেয়াল
তুলছে।

%%K t05-01-16 p
তার কাছে একটি করণী, একটি গামলা ও একটি \VUgras{ওলন} \textsuperscript{(109)}
আছে। কিন্তু এই নামগুলি তোমার মনে থাকবে না।

%%K t05-01-17 p
\textsc{জন}। --- থাকবে না কেন, আমি কাকার কাছে একটি ছোট করণী ও একটি গামলা
চাইব, আর ওলন তো আমি সুতো দিয়ে নিজেই বানিয়ে নেব।

%%K t05-tit-3 sub
\VUsaut{3.0mm}
\VUpk{10.2pt}{40}{উপকরণ।}
\VUsaut{3.0mm}

%%K t05-01-18 p
\textsc{কাকা}। --- বাহ, খুব ভালো। কিন্তু বলো তো জন, বাড়ি বানাতে কী কী দরকার?

%%K t05-01-19 p
\textsc{জন}। --- \VUgras{কাটা পাথর} \textsuperscript{(59)} ও \VUgras{মশলা}
\textsuperscript{(65)} দরকার।

%%K t05-01-20 p
\textsc{কাকা}। --- ঠিক। ওই বড় টুপিওয়ালা লোকটিকে দেখো, নিজের \VUgras{গাড়িতে}
\textsuperscript{(136)}। সে \VUgras{গাড়োয়ান} \textsuperscript{(135)}। সে প্রতিদিন
এখানে \VUgras{পাথরের চাঁই} \textsuperscript{(60)} আনে। সে নিজের গাড়ি উঠানে
খালি করে, আর \VUgras{পাথরকাটারা} \textsuperscript{(83)} চাঁইগুলি কেটে নিজেদের
\VUgras{পাথরের করাত} \textsuperscript{(85)} দিয়ে চেরে। একজন \VUgras{জোগানদার}
\textsuperscript{(146)} সেগুলি রাজমিস্ত্রির কাছে পৌঁছে দেয়, যে সেগুলি নিজের
জায়গায় বসায়।

%%K t05-01-21 p
এদিকে \VUgras{মশলা-মাখানেওয়ালা} \textsuperscript{(87)} মশলা তৈরি করে। তার
\VUgras{বালি} \textsuperscript{(62)}, \VUgras{চুন} \textsuperscript{(63)} ও
\VUgras{জল} \textsuperscript{(64)} দরকার। চুন তার পাশে একটি \VUgras{বস্তায়}
\textsuperscript{(71)} আছে; একজন \VUgras{রাজ-সহায়ক} \textsuperscript{(111)} তাকে
\VUgras{ঠেলাগাড়িতে} \textsuperscript{(112)} বালি এনে দেয়, আর
\VUgras{শিক্ষানবিশ} \textsuperscript{(113)} পাম্পের \textsuperscript{(58)} কাছে।
সে একটি \VUgras{বালতি} \textsuperscript{(114)} ভরছে। মশলা শিগগিরই তৈরি হয়ে
যাবে। \VUgras{মশলা-বাহক} \textsuperscript{(107)} সেটি নিয়ে সিঁড়ি বেয়ে ভারায়
তুলে নিয়ে যায়, যেখানে একজন রাজমিস্ত্রি বাড়ির সেই দিকটি শেষ করছে।

%%K t05-01-22 p
আর এখন বলো তো, যে মজুর \VUgras{ছাদ} \textsuperscript{(31)} বানায় তাকে কী বলে?

%%K t05-01-23 p
\textsc{জন}। --- সে \VUgras{ছাদ-নির্মাতা} \textsuperscript{(118)}।

%%K t05-tit-4 sub
\VUsaut{3.0mm}
\VUpk{10.2pt}{40}{বাড়ির ছাদ।}
\VUsaut{3.0mm}

%%K t05-01-24 p
\textsc{শ্রী শ্বেত}। --- হ্যাঁ, কিন্তু আগে আসে \VUgras{ছুতোর}
\textsuperscript{(115)}।

%%K t05-01-25 p
ছুতোর \VUgras{ছাদের কাঠামো} \textsuperscript{(35)} বানায়। সে \VUgras{কড়িগুলিকে}
\textsuperscript{(37)} \VUgras{গোঁজ} \textsuperscript{(117)} দিয়ে জোড়ে। উপরে ওই
মজুরেরা, নিজেদের চওড়া মখমলি নিকারে, ছুতোর। একজন একটি ছোট কড়ি ধরে আছে আর
অন্যজন নিজের \VUgras{কুড়ুল} \textsuperscript{(116)} দিয়ে একটি গোঁজ গড়ছে।
\VUgras{চিমনির} \textsuperscript{(41)} পাশের জন ছাদ-নির্মাতা। সে
\VUgras{বরগার} (*) উপরে পাটাতন ঠোকে, আর পাটাতনের উপরে \VUgras{স্লেট}
\textsuperscript{(66)}। কখনও কখনও সে খাপরা বিছায়।

%%K t05-noto-1 noto
\VUnotes{143.2mm}{%
(*) \VUgras{Balk-o} = জার্মান \textit{Balken}, ইংরেজি \textit{joist},
ফরাসি \textit{solive}, ইতালীয় \textit{travicello}, রুশ \textit{balka},
স্প্যানিশ ও পর্তুগিজ \textit{viga}। একটি কারিগরি ধাতু যা এখনও গৃহীত হয়নি,
তবু এখানে ফরাসি শব্দ \textit{solive}-এর অর্থ দিতে প্রস্তাবিত, যা
\textit{trabo} (ফরাসি \textit{poutre}) নয়, ছোট কড়িও নয়। সেটি চৌকো করে
গড়া কাঠ, যা একটি মেঝে ও একটি ছাদ ধরে রাখে।%
}

%%K t05-01-26 p
আস্তাবলের ছাদ \VUgras{খাপরার} \textsuperscript{(55)}, বাড়ির
\VUgras{স্লেটের} \textsuperscript{(32)}, আর বারান্দার \VUgras{দস্তার}
\textsuperscript{(24)}।

%%K t05-01-27 p
\textsc{জন}। --- দস্তার ছাদও কি ছাদ-নির্মাতাই বানায়?

%%K t05-01-28 p
\textsc{শ্রী শ্বেত}। --- না, সে \VUgras{দস্তাকার} \textsuperscript{(121)} বা
\VUgras{নলকারের} কাজ --- সেই লোকই যে বাড়ির ছাদে একটি
\VUgras{ঘুলঘুলি-জানালা} \textsuperscript{(38)} বসাচ্ছে। সে-ই
\VUgras{ছাদের নালাও} \textsuperscript{(43)} ও \VUgras{নামার নল}
\textsuperscript{(42)} লাগায়। সে দস্তা ও টিন ঝালাই করে। সে-ই
\VUgras{বজ্রনিরোধক} \textsuperscript{(40)} ও \VUgras{বাতাস-সূচক}
\textsuperscript{(39)} \VUgras{কার্নিশে} \textsuperscript{(36)} বসিয়েছিল।

%%K t05-01-29 p
\textsc{জন}। --- তা হলে কি বাড়ি শেষ হয়ে গেল?

%%K t05-tit-5 sub
\VUsaut{3.0mm}
\VUpk{10.2pt}{40}{যন্ত্রপাতি।}
\VUsaut{3.0mm}

%%K t05-01-30 p
\textsc{শ্রী শ্বেত}। --- শেষ হয়নি। \VUgras{পলেস্তারাকার}
\textsuperscript{(99)} \VUgras{ইট} \textsuperscript{(61)} দিয়ে সেই দেয়ালগুলি
তোলে যা বাড়িটিকে আলাদা আলাদা ঘরে ভাগ করে। সে \VUgras{পলেস্তারা}
\textsuperscript{(70)} \VUgras{ছাদে} \textsuperscript{(26)} ও দেয়ালে চড়ায়।
এখন সে \VUgras{বারান্দার} \textsuperscript{(33)} ছাদে কাজ করছে। রাজমিস্ত্রির
মতো তার কাছেও একটি \VUgras{করণী} \textsuperscript{(100)} ও একটি \VUgras{গামলা}
\textsuperscript{(101)} আছে।

%%K t05-01-31 p
তারপর ছুতোর দরজা, জানালা, \VUgras{কাঠের মেঝে} \textsuperscript{(16)} ও
\VUgras{ঝিলিমিলি} \textsuperscript{(19)} বানায়।

%%K t05-01-32 p
এখন সে উঠানে নিজের \VUgras{কাজের টেবিলে} \textsuperscript{(90)} কাজ করছে। তার
কাছে একটি \VUgras{করাত} \textsuperscript{(94)} আছে, \VUgras{কাঠ}
\textsuperscript{(69)} চেরার জন্য, একটি \VUgras{রান্দা} \textsuperscript{(91)},
একটি \VUgras{কষা} \textsuperscript{(92)}, একটি \VUgras{বাটালি}
\textsuperscript{(94 \textit{bis})}, একটি \VUgras{কম্পাস} \textsuperscript{(95)} ও
কাঠের একটি \VUgras{মুগুর} \textsuperscript{(95 \textit{bis})}।

%%K t05-01-33 p
\textsc{জন}। --- আরে, তা হলে তো ছুতোর হওয়া রাজমিস্ত্রি হওয়ার চেয়েও বেশি
মজার। কাকা, আপনি আমাকে একটি কাজের টেবিল, একটি করাত ও একটি রান্দা কিনে দেবেন?
আমি মেদরের জন্য একটি ঘর বানাব।

%%K t05-01-34 p
\textsc{কাকা}। --- হ্যাঁ বাছা।

%%K t05-01-35 p
\textsc{জন}। --- কী আনন্দ হল! আর ওই লোকটি কে, যে সামনের দরজার পাশে
দাঁড়িয়ে?

%%K t05-tit-6 sub
\VUsaut{3.0mm}
\VUpk{10.2pt}{40}{রং করা।}
\VUsaut{3.0mm}

%%K t05-01-36 p
\textsc{শ্রী শ্বেত}। --- সে \VUgras{রঙমিস্ত্রি} \textsuperscript{(102)}। সে
নিজের মইয়ে দাঁড়িয়ে, \VUgras{রঙের} \textsuperscript{(103)} একটি টিন ও নিজের
\VUgras{তুলি} \textsuperscript{(104)} নিয়ে। সে সামনের দরজার কাঠে রং চড়াচ্ছে
যাতে তা বেশি দিন টেকে।

%%K t05-01-37 p
ঘরের দেয়ালের কাগজও সে-ই সাঁটে।

%%K t05-01-38 p
\VUgras{সাজসজ্জাকার} \textsuperscript{(107)}, যে \VUgras{মইয়ে}
\textsuperscript{(106)} উঠে \VUgras{বারান্দায়} \textsuperscript{(28)} আছে, একটি
\VUgras{পরদা} \textsuperscript{(29)} লাগাচ্ছে।

%%K t05-01-39 p
বড় জানালার পাশে দাঁড়ানো মজুরটি \VUgras{কাচমিস্ত্রি} \textsuperscript{(96)}।
সে \VUgras{কাচ} \textsuperscript{(18)} \VUgras{মশলা} \textsuperscript{(73)}
দিয়ে জমায়। ওই অন্য মজুরটি, যে ভাঁড়ারঘরের দরজার পাশে হাঁটু গেড়ে বসে,
\VUgras{তালামিস্ত্রি} \textsuperscript{(97)}। সে একটি \VUgras{তালা}
\textsuperscript{(6)} বসাচ্ছে। তার \VUgras{রেতি} \textsuperscript{(98)} তার পাশে
পড়ে আছে।

%%K t05-01-40 p
\textsc{জন}। --- যে লোকটি নিজের \VUgras{হাতুড়ি} \textsuperscript{(84)} দিয়ে
লোহার টুকরোয় ঘা দিচ্ছে, সেও কি তালামিস্ত্রি?

%%K t05-01-41 p
\textsc{শ্রী শ্বেত}। --- ঠিক বলতে গেলে সে \VUgras{কামার}
\textsuperscript{(125)}। সে \VUgras{লোহার জিনিস} \textsuperscript{(67)} গড়ছে
\VUgras{বিদ্যুৎকর্মীর} \textsuperscript{(124)} জন্য, যে \VUgras{গাড়িঘরের}
\textsuperscript{(51)} ছাদে আছে। তার কাছে একটি \VUgras{হাপর}
\textsuperscript{(126)}, একটি \VUgras{নেহাই} \textsuperscript{(127)} ও একটি
\VUgras{সাঁড়াশি} \textsuperscript{(128)} আছে।

%%K t05-01-42 p
বিদ্যুৎকর্মী টেলিফোনের \VUgras{তামার তার} \textsuperscript{(44)} জুড়ছে।

%%K t05-tit-7 sub
\VUpk{10.2pt}{40}{বাগান করা।}
\VUsaut{3.0mm}

%%K t05-01-43 p
\textsc{কাকা}। --- \VUgras{উঠানের} \textsuperscript{(45)} শেষে,
\VUgras{জাফরির} \textsuperscript{(47)} ও \VUgras{গাড়ির ফটকের}
\textsuperscript{(48)} পাশে, তোমার চোখে পড়ছে \VUgras{মালি}
\textsuperscript{(129)}। সে ছোট \VUgras{বাগানটি} \textsuperscript{(49)} তৈরি
করছে। সে নিজের \VUgras{কোদাল} \textsuperscript{(131)} দিয়ে মাটি খুঁড়ে
ফেলেছে, বীজ বুনছে ও \VUgras{আঁচড়া} \textsuperscript{(130)} দিয়ে সমান করছে।
তারপর নিজের \VUgras{ঝারি} \textsuperscript{(132)} দিয়ে সে \VUgras{কেয়ারিগুলিতে}
\textsuperscript{(150)} জল দেবে আর গাছ গজাবে।

%%K t05-01-44 p
\VUgras{সহিস} \textsuperscript{(133)}, যে এইমাত্র ঘোড়ার দেখাশোনা করে তাকে
খাইয়েছে, \VUgras{গাড়িটি} \textsuperscript{(52)} সাফ করছে --- স্পঞ্জ, একটি
\VUgras{হাতপাম্প} \textsuperscript{(134)} ও জলের বালতি দিয়ে। তারপর সে সেটি
গাড়িঘরে রেখে দেবে আর ভারী \VUgras{খিল} \textsuperscript{(53)} দিয়ে দরজা বন্ধ
করবে।

%%K t05-01-45 p
এই তো, এবার নিশ্চয় সন্তুষ্ট হলে; ব্যাখ্যা যথেষ্ট হয়েছে।

%%K t05-01-46 p
\textsc{জন}। --- হ্যাঁ কাকা, কিন্তু আমি বাড়ির ভিতরটাও দেখতে চাই।

%%K t05-01-47 p
\textsc{কাকা}। --- সেটি কালকের জন্য। শ্রী রু তোমাকে নকশা বুঝিয়ে দেবেন ও
প্রতিটি ঘরের নাম বলবেন।

%%K t05-tit-8 sub
\VUsaut{3.0mm}
\VUpk{10.2pt}{40}{বাড়ির ভিতরের বিন্যাস।}
\VUsaut{2.4mm}

%%K t05-apar-2 apar
{\centering\textit{(নকশা দেখুন।)}\par}
\VUsaut{2.4mm}

%%K t05-01-48 p
\textsc{স্থপতি}। --- এসো জন, আমি তোমাকে বাড়ির আলাদা আলাদা অংশ দেখাই।

%%K t05-01-49 p
চলো \VUgras{নিচের তলায়} \textsuperscript{(7)} যাই। এই যে \VUgras{ঘণ্টির} বোতাম
\textsuperscript{(11)}। আমরা তিনটি \VUgras{ধাপ} \textsuperscript{(14)} উঠি, যা
\VUgras{সিঁড়ি-চাতালের} \textsuperscript{(9)}, তারপর \VUgras{দোরগোড়া}
\textsuperscript{(10)} পেরোই, যা \VUgras{সামনের দরজার} \textsuperscript{(8)}, আর
এই যে আমরা \VUgras{সিঁড়ির} \textsuperscript{(13)} নিচে।

%%K t05-01-50 p
\textsc{জন}। --- এটি বারান্দা-পথ।

%%K t05-01-51 p
\textsc{স্থপতি}। --- না, এটি \VUgras{প্রবেশকক্ষ} \textsuperscript{(12)}।
\VUgras{বারান্দা-পথ} \textsuperscript{(a)} শেষ দিকে। ডান দিকে এই যে
\VUgras{বৈঠকখানা} \textsuperscript{(b)} ও \VUgras{খাবারঘর} \textsuperscript{(c)};
খাবারঘরের সামনে \VUgras{রান্নাঘর} \textsuperscript{(d)} ও \VUgras{ভাঁড়ার}
\textsuperscript{(e)}; তারপর সামনের দিকে \VUgras{পড়ার ঘর} \textsuperscript{(f)}।
\VUgras{বারান্দা} \textsuperscript{(23)}, তার \VUgras{কাচের দরজাসহ}
\textsuperscript{(25)}, বৈঠকখানার পাশে।

%%K t05-01-52 p
এবার আমরা সিঁড়ি উঠব। \VUgras{রেলিং} \textsuperscript{(15)} ধরে থাকো আর ধাপ
গোনো। চোদ্দোটি। এই যে \VUgras{চাতাল} \textsuperscript{(m)}: আমরা প্রথম
\VUgras{তলায়} \textsuperscript{(27)}।

%%K t05-tit-9 sub
\VUsaut{3.0mm}
\VUpk{10.2pt}{40}{প্রথম তলা।}
\VUsaut{3.0mm}

%%K t05-01-53 p
সিঁড়ির সামনে \VUgras{শৌচাগার} \textsuperscript{(20)} ও \VUgras{সাজঘর}
\textsuperscript{(j)}। ডান দিকে একটি সুন্দর \VUgras{শোবার ঘর}
\textsuperscript{(g)}, যার দুটি জানালা বাড়ির \VUgras{সামনের দিকে}
\textsuperscript{(1)} খোলে। বাঁ দিকে \VUgras{স্নানঘর} \textsuperscript{(1)} ও
\VUgras{অতিথিকক্ষ} \textsuperscript{(i)}, যেখানে একটি বড় \VUgras{কুলুঙ্গি}
\textsuperscript{(h)} আছে। শোবার ঘরের পাশে \VUgras{ছেলেমেয়েদের ঘর}
\textsuperscript{(k)}। সেটি একটি চওড়া \VUgras{ছাদ-বারান্দায়}
\textsuperscript{(21)} খোলে। সেই বারান্দার নিচে আর একটি শৌচাগার
\textsuperscript{(20)} আছে, আর দেয়াল ঘেঁষে এই যে \VUgras{ঘণ্টা}
\textsuperscript{(22)}, যা খাওয়ার জন্য বাজানো হয়।

%%K t05-01-54 p
\textsc{জন}। --- চলুন \VUgras{দ্বিতীয় তলায়} যাই।

%%K t05-01-55 p
\textsc{স্থপতি}। --- দ্বিতীয় তলা তো নেই-ই। ছাদের নিচে কেবল
\VUgras{চিলেকোঠা} \textsuperscript{(33)} ও গুদাম \textsuperscript{(34)}। আমাদের
ঘোরা শেষ, এবার নামা যায়।

%%K t05-01-56 p
\textsc{জন}। --- ধন্যবাদ মহাশয়, খুব কৌতূহলজনক লাগল। যা দেখলাম সব আমি বাবাকে
বলব।
\VUsaut{4.0mm}
\VUfilet{22mm}
